\sectionheading{5.42 Loop Pattern}

Loop ব্যবহার করে কোনো নির্দিষ্ট বিন্যাসে Character, Symbol বা Number
বারবার Print করাকে সাধারণভাবে \textbf{Loop Pattern} বলা হয়।
Pattern তৈরিতে সাধারণত \textbf{Nested Loop} বেশি ব্যবহৃত হয়।

%================================================
% Basic Idea
%================================================

\booksubheading{Pattern তৈরির মূল ধারণা}

Pattern-কে সাধারণত \textbf{Row} ও \textbf{Column} হিসেবে চিন্তা করা হয়।

\begin{rememberbox}

\begin{center}

\textbf{Outer Loop}
\(\longrightarrow\)
\textbf{Row নিয়ন্ত্রণ করে}

\qquad

\textbf{Inner Loop}
\(\longrightarrow\)
\textbf{প্রতি Row-এর Item / Column নিয়ন্ত্রণ করে}

\medskip

\textbf{Outer Loop = কতটি Row}
\qquad
\textbf{Inner Loop = প্রতি Row-তে কতটি Item}

\end{center}

\end{rememberbox}

%================================================
% print and end
%================================================

\booksubheading{\textenglish{\ttfamily print()} ও \textenglish{\ttfamily end} কেন গুরুত্বপূর্ণ?}

\begin{codeoutputbox}{সাধারণ \textenglish{\ttfamily print()}}

\noindent
\begin{minipage}[t]{0.53\linewidth}
\codeheading

\begin{english}
\ttfamily
print("*")\\
print("*")
\end{english}

\end{minipage}
\hfill
\begin{minipage}[t]{0.43\linewidth}
\outputheading

\begin{english}
\ttfamily
*\\
*
\end{english}

প্রতিটি
\textenglish{\ttfamily print()}
শেষে নতুন Line তৈরি হয়।

\end{minipage}

\end{codeoutputbox}

\begin{codeoutputbox}{\textenglish{\ttfamily end=" "} ব্যবহার}

\noindent
\begin{minipage}[t]{0.53\linewidth}
\codeheading

\begin{english}
\ttfamily
print("*", end=" ")\\
print("*", end=" ")
\end{english}

\end{minipage}
\hfill
\begin{minipage}[t]{0.43\linewidth}
\outputheading

\textenglish{\ttfamily * *}

\medskip

\textenglish{\ttfamily end=" "}
ব্যবহার করলে পরবর্তী Output একই Line-এ থাকে।

\end{minipage}

\end{codeoutputbox}

\begin{rememberbox}

Pattern-এর সাধারণ Structure—

\begin{codebox}
\begin{english}
\ttfamily
for ...:\\
\hspace*{4mm}for ...:\\
\hspace*{8mm}print(..., end=" ")\\
\hspace*{4mm}print()
\end{english}
\end{codebox}

ভিতরের
\textenglish{\ttfamily print(..., end=" ")}
একই Row-তে Item Print করে এবং শেষের
\textenglish{\ttfamily print()}
নতুন Row শুরু করে।

\end{rememberbox}

%================================================
% Rectangular Pattern
%================================================

\booksubheading{১. Rectangular Star Pattern}

\begin{codeoutputbox}{Rectangular Pattern}

\noindent
\begin{minipage}[t]{0.53\linewidth}
\codeheading

\begin{english}
\ttfamily
for i in range(3):\\
\hspace*{4mm}for j in range(4):\\
\hspace*{8mm}print("*", end=" ")\\
\hspace*{4mm}print()
\end{english}

\end{minipage}
\hfill
\begin{minipage}[t]{0.43\linewidth}
\outputheading

\begin{english}
\ttfamily
* * * *\\
* * * *\\
* * * *
\end{english}

\medskip

Outer Loop = 3 Row

Inner Loop = প্রতি Row-তে 4 Star

\medskip

মোট Star =
\(3 \times 4 = 12\)

\end{minipage}

\end{codeoutputbox}

\begin{tabularx}
{\linewidth}
{
|>{\bfseries\centering\arraybackslash}p{0.25\linewidth}
|>{\centering\arraybackslash}Y
|>{\centering\arraybackslash}Y|
}

\hline
\rowcolor{softBlue}

\textbf{Row}
&
\textbf{Inner Loop Iteration}
&
\textbf{Star}
\\
\hline

1 & 4 & 4 \\
\hline
2 & 4 & 4 \\
\hline
3 & 4 & 4 \\
\hline

\end{tabularx}

\begin{rememberbox}

\textbf{\color{deepOrange}MCQ Shortcut:}
যদি Outer Loop
\textenglish{\ttfamily m}
বার এবং Inner Loop প্রতিবার
\textenglish{\ttfamily n}
বার চলে, তবে মোট Inner Execution =
\(m \times n\)।

\end{rememberbox}

%================================================
% Increasing Pattern
%================================================

\booksubheading{২. Increasing Star Pattern}

\begin{codeoutputbox}{Increasing Pattern}

\noindent
\begin{minipage}[t]{0.53\linewidth}
\codeheading

\begin{english}
\ttfamily
for i in range(1, 5):\\
\hspace*{4mm}for j in range(i):\\
\hspace*{8mm}print("*", end=" ")\\
\hspace*{4mm}print()
\end{english}

\end{minipage}
\hfill
\begin{minipage}[t]{0.43\linewidth}
\outputheading

\begin{english}
\ttfamily
*\\
* *\\
* * *\\
* * * *
\end{english}

\medskip

\textenglish{\ttfamily i = 1,2,3,4}

তাই
\textenglish{\ttfamily range(i)}
যথাক্রমে 1, 2, 3 ও 4 বার চলে।

\end{minipage}

\end{codeoutputbox}

\begin{rememberbox}

Increasing Pattern-এর সাধারণ Logic—

\begin{codebox}
\begin{english}
\ttfamily
for i in range(1, n + 1):\\
\hspace*{4mm}for j in range(i):\\
\hspace*{8mm}print("*", end=" ")\\
\hspace*{4mm}print()
\end{english}
\end{codebox}

\textbf{Increasing Pattern}
\(\longrightarrow\)
Row Number বাড়ার সঙ্গে Item সংখ্যাও বাড়ে।

\end{rememberbox}

%================================================
% Decreasing Pattern
%================================================

\booksubheading{৩. Decreasing Star Pattern}

\begin{codeoutputbox}{Decreasing Pattern}

\noindent
\begin{minipage}[t]{0.53\linewidth}
\codeheading

\begin{english}
\ttfamily
for i in range(4, 0, -1):\\
\hspace*{4mm}for j in range(i):\\
\hspace*{8mm}print("*", end=" ")\\
\hspace*{4mm}print()
\end{english}

\end{minipage}
\hfill
\begin{minipage}[t]{0.43\linewidth}
\outputheading

\begin{english}
\ttfamily
* * * *\\
* * *\\
* *\\
*
\end{english}

\medskip

\textenglish{\ttfamily i}
এর Value:
\textenglish{\ttfamily 4,3,2,1}

তাই প্রতি Row-তে Star সংখ্যা কমে।

\end{minipage}

\end{codeoutputbox}

%================================================
% Alternative decreasing
%================================================

\booksubheading{Decreasing Pattern-এর আরেকটি Logic}

\begin{codeoutputbox}{Formula ব্যবহার করে}

\noindent
\begin{minipage}[t]{0.56\linewidth}
\codeheading

\begin{english}
\ttfamily
n = 4\\
for i in range(1, n + 1):\\
\hspace*{4mm}for j in range(n - i + 1):\\
\hspace*{8mm}print("*", end=" ")\\
\hspace*{4mm}print()
\end{english}

\end{minipage}
\hfill
\begin{minipage}[t]{0.40\linewidth}
\outputheading

প্রতি Row-এর Item সংখ্যা—

\[
n-i+1
\]

যদি
\textenglish{\ttfamily n = 4},
তবে Item সংখ্যা—

\begin{center}
4, 3, 2, 1
\end{center}

তাই Output আগের Pattern-এর মতোই।

\end{minipage}

\end{codeoutputbox}

%================================================
% Row Number Pattern
%================================================

\booksubheading{৪. Row Number Pattern}

\begin{codeoutputbox}{Row Number Pattern}

\noindent
\begin{minipage}[t]{0.53\linewidth}
\codeheading

\begin{english}
\ttfamily
for i in range(1, 5):\\
\hspace*{4mm}for j in range(i):\\
\hspace*{8mm}print(i, end=" ")\\
\hspace*{4mm}print()
\end{english}

\end{minipage}
\hfill
\begin{minipage}[t]{0.43\linewidth}
\outputheading

\begin{english}
\ttfamily
1\\
2 2\\
3 3 3\\
4 4 4 4
\end{english}

\medskip

Print করা হচ্ছে
\textenglish{\ttfamily i}।

তাই প্রতিটি Row-তে Row Number দেখা যায়।

\end{minipage}

\end{codeoutputbox}

%================================================
% Column Number Pattern
%================================================

\booksubheading{৫. Column Number Pattern}

\begin{codeoutputbox}{Column Number Pattern}

\noindent
\begin{minipage}[t]{0.53\linewidth}
\codeheading

\begin{english}
\ttfamily
for i in range(1, 5):\\
\hspace*{4mm}for j in range(1, i + 1):\\
\hspace*{8mm}print(j, end=" ")\\
\hspace*{4mm}print()
\end{english}

\end{minipage}
\hfill
\begin{minipage}[t]{0.43\linewidth}
\outputheading

\begin{english}
\ttfamily
1\\
1 2\\
1 2 3\\
1 2 3 4
\end{english}

\medskip

Print করা হচ্ছে
\textenglish{\ttfamily j}।

তাই Inner Loop-এর Value Output-এ দেখা যায়।

\end{minipage}

\end{codeoutputbox}

\begin{rememberbox}

\begin{center}

\textenglish{\ttfamily print(i, end=" ")}
\(\longrightarrow\)
\textbf{Row Value}

\qquad

\textenglish{\ttfamily print(j, end=" ")}
\(\longrightarrow\)
\textbf{Column Value}

\end{center}

\end{rememberbox}

%================================================
% i vs j
%================================================

\booksubheading{\textenglish{\ttfamily i} বনাম \textenglish{\ttfamily j}}

\begin{tabularx}
{\linewidth}
{
|>{\bfseries\centering\arraybackslash}p{0.25\linewidth}
|Y|
}

\hline
\rowcolor{softBlue}

\textbf{Variable}
&
\textbf{সাধারণ কাজ}
\\
\hline

\textenglish{\ttfamily i}
&
Outer Loop / Row নির্দেশ করতে ব্যবহৃত হয়।
\\
\hline

\textenglish{\ttfamily j}
&
Inner Loop / Column বা Row-এর Item নির্দেশ করতে ব্যবহৃত হয়।
\\
\hline

\end{tabularx}

\begin{rememberbox}

\textenglish{\ttfamily i}
ও
\textenglish{\ttfamily j}
বাধ্যতামূলক Variable Name নয়।
অন্য বৈধ Variable Name-ও ব্যবহার করা যায়।

\end{rememberbox}

%================================================
% Square Number Pattern
%================================================

\booksubheading{৬. Square Number Pattern}

\begin{codeoutputbox}{Square Number Pattern}

\noindent
\begin{minipage}[t]{0.53\linewidth}
\codeheading

\begin{english}
\ttfamily
for i in range(1, 4):\\
\hspace*{4mm}for j in range(1, 4):\\
\hspace*{8mm}print(j, end=" ")\\
\hspace*{4mm}print()
\end{english}

\end{minipage}
\hfill
\begin{minipage}[t]{0.43\linewidth}
\outputheading

\begin{english}
\ttfamily
1 2 3\\
1 2 3\\
1 2 3
\end{english}

\medskip

Outer Loop = 3 Row

Inner Loop প্রতিবার
\textenglish{\ttfamily 1,2,3}
Print করে।

\end{minipage}

\end{codeoutputbox}

%================================================
% Same Number Square
%================================================

\booksubheading{৭. একই Number-এর Square Pattern}

\begin{codeoutputbox}{একই Row Number বারবার}

\noindent
\begin{minipage}[t]{0.53\linewidth}
\codeheading

\begin{english}
\ttfamily
for i in range(1, 4):\\
\hspace*{4mm}for j in range(3):\\
\hspace*{8mm}print(i, end=" ")\\
\hspace*{4mm}print()
\end{english}

\end{minipage}
\hfill
\begin{minipage}[t]{0.43\linewidth}
\outputheading

\begin{english}
\ttfamily
1 1 1\\
2 2 2\\
3 3 3
\end{english}

\medskip

Inner Loop কতবার Print হবে তা নিয়ন্ত্রণ করে।

Print হচ্ছে Outer Loop-এর
\textenglish{\ttfamily i}।

\end{minipage}

\end{codeoutputbox}

%================================================
% Continuous Number Pattern
%================================================

\booksubheading{৮. Continuous Number Pattern}

\begin{codeoutputbox}{Continuous Number Pattern}

\noindent
\begin{minipage}[t]{0.55\linewidth}
\codeheading

\begin{english}
\ttfamily
num = 1\\
for i in range(1, 4):\\
\hspace*{4mm}for j in range(i):\\
\hspace*{8mm}print(num, end=" ")\\
\hspace*{8mm}num += 1\\
\hspace*{4mm}print()
\end{english}

\end{minipage}
\hfill
\begin{minipage}[t]{0.41\linewidth}
\outputheading

\begin{english}
\ttfamily
1\\
2 3\\
4 5 6
\end{english}

\medskip

\textenglish{\ttfamily num}
একটি আলাদা Counter।

প্রতিবার
\textenglish{\ttfamily num += 1}
হওয়ায় Number ধারাবাহিকভাবে বাড়ে।

\end{minipage}

\end{codeoutputbox}

%================================================
% Pattern solving technique
%================================================

\booksubheading{Pattern তৈরির সহজ কৌশল}

Pattern দেখলে প্রথমে তিনটি প্রশ্ন করো—

\begin{tabularx}
{\linewidth}
{
|>{\bfseries\RaggedRight\arraybackslash}p{0.34\linewidth}
|Y|
}

\hline
\rowcolor{softBlue}

\textbf{প্রশ্ন}
&
\textbf{কী নির্ধারণ করবে?}
\\
\hline

মোট কতটি Row?
&
Outer Loop।
\\
\hline

প্রতি Row-তে কতটি Item?
&
Inner Loop।
\\
\hline

কী Print করতে হবে?
&
\textenglish{\ttfamily "*"},
\textenglish{\ttfamily i},
\textenglish{\ttfamily j}
বা আলাদা Counter।
\\
\hline

\end{tabularx}

\begin{rememberbox}

\begin{center}

\textbf{কত Row?}
\(\longrightarrow\)
\textbf{Outer Loop}

\qquad

\textbf{কত Item?}
\(\longrightarrow\)
\textbf{Inner Loop}

\medskip

\textbf{কী Print হবে?}
\(\longrightarrow\)
\textenglish{\ttfamily * / i / j / counter}

\end{center}

\end{rememberbox}

%================================================
% Increasing Analysis
%================================================

\booksubheading{Increasing Pattern বিশ্লেষণের Shortcut}

\begin{tabularx}
{\linewidth}
{
|>{\bfseries\centering\arraybackslash}p{0.25\linewidth}
|>{\centering\arraybackslash}p{0.25\linewidth}
|Y|
}

\hline
\rowcolor{softBlue}

\textbf{Row \textenglish{\ttfamily i}}
&
\textbf{Star সংখ্যা}
&
\textbf{Logic}
\\
\hline

1 & 1 & \textenglish{\ttfamily range(1)} \\
\hline
2 & 2 & \textenglish{\ttfamily range(2)} \\
\hline
3 & 3 & \textenglish{\ttfamily range(3)} \\
\hline
4 & 4 & \textenglish{\ttfamily range(4)} \\
\hline

\end{tabularx}

\begin{center}

\textbf{Star সংখ্যা = Row Number}
\(\Rightarrow\)
Inner Loop =
\textenglish{\ttfamily range(i)}

\end{center}

%================================================
% Decreasing Analysis
%================================================

\booksubheading{Decreasing Pattern বিশ্লেষণের Shortcut}

যদি মোট Row
\textenglish{\ttfamily n}
হয়—

\begin{tabularx}
{\linewidth}
{
|>{\bfseries\centering\arraybackslash}p{0.35\linewidth}
|>{\centering\arraybackslash}Y|
}

\hline
\rowcolor{softBlue}

\textbf{Row \textenglish{\ttfamily i}}
&
\textbf{Item সংখ্যা}
\\
\hline

1 & \textenglish{\ttfamily n} \\
\hline

2 & \textenglish{\ttfamily n-1} \\
\hline

3 & \textenglish{\ttfamily n-2} \\
\hline

\(\cdots\) & \(\cdots\) \\
\hline

\end{tabularx}

\begin{center}

\textbf{সাধারণ Formula:}
\qquad
\(n-i+1\)

\end{center}

%================================================
% Nested Loop Working
%================================================

\booksubheading{Nested Loop কীভাবে চলে?}

\begin{codeoutputbox}{Nested Loop Execution}

\noindent
\begin{minipage}[t]{0.52\linewidth}
\codeheading

\begin{english}
\ttfamily
for i in range(2):\\
\hspace*{4mm}for j in range(3):\\
\hspace*{8mm}print(i, j)
\end{english}

\end{minipage}
\hfill
\begin{minipage}[t]{0.44\linewidth}
\outputheading

\begin{english}
\ttfamily
0 0\\
0 1\\
0 2\\
1 0\\
1 1\\
1 2
\end{english}

\medskip

\textenglish{\ttfamily i = 0}
হলে
\textenglish{\ttfamily j = 0,1,2}।

তারপর
\textenglish{\ttfamily i = 1}
হলেও Inner Loop আবার সম্পূর্ণ চলে।

\end{minipage}

\end{codeoutputbox}

\begin{rememberbox}

\begin{center}

\textbf{Outer Loop-এর প্রতিটি Iteration-এর জন্য
Inner Loop সম্পূর্ণভাবে চলে।}

\end{center}

\end{rememberbox}

%================================================
% print position
%================================================

\booksubheading{Pattern-এ \textenglish{\ttfamily print()}-এর অবস্থান}

\begin{codeoutputbox}{সঠিক Position}

\noindent
\begin{minipage}[t]{0.53\linewidth}
\codeheading

\begin{english}
\ttfamily
for i in range(3):\\
\hspace*{4mm}for j in range(3):\\
\hspace*{8mm}print("*", end=" ")\\
\hspace*{4mm}print()
\end{english}

\end{minipage}
\hfill
\begin{minipage}[t]{0.43\linewidth}
\outputheading

শেষের
\textenglish{\ttfamily print()}
Inner Loop-এর বাইরে কিন্তু Outer Loop-এর ভিতরে।

তাই একটি Row-এর সব Star Print হওয়ার পর নতুন Line হয়।

\end{minipage}

\end{codeoutputbox}

\begin{codeoutputbox}{ভুল Position}

\noindent
\begin{minipage}[t]{0.53\linewidth}
\codeheading

\begin{english}
\ttfamily
for i in range(3):\\
\hspace*{4mm}for j in range(3):\\
\hspace*{8mm}print("*", end=" ")\\
\hspace*{8mm}print()
\end{english}

\end{minipage}
\hfill
\begin{minipage}[t]{0.43\linewidth}
\outputheading

দ্বিতীয়
\textenglish{\ttfamily print()}
Inner Loop-এর ভিতরে।

তাই প্রতিটি Star-এর পর নতুন Line তৈরি হবে এবং প্রত্যাশিত Row Pattern হবে না।

\end{minipage}

\end{codeoutputbox}

\begin{rememberbox}

Row পরিবর্তনের খালি
\textenglish{\ttfamily print()}
সাধারণত—

\begin{center}

\textbf{Inner Loop-এর বাইরে}
\qquad
কিন্তু
\qquad
\textbf{Outer Loop-এর ভিতরে}

\end{center}

থাকে।

\end{rememberbox}

%================================================
% end="" vs end=" "
%================================================

\booksubheading{\textenglish{\ttfamily end=""} বনাম \textenglish{\ttfamily end=" "}}

\begin{codeoutputbox}{Output Formatting}

\noindent
\begin{minipage}[t]{0.48\linewidth}

\textbf{\color{darkTeal}Space ছাড়া}

\begin{english}
\ttfamily
for i in range(4):\\
\hspace*{4mm}print("*", end="")
\end{english}

\medskip

\textbf{Output:}
\textenglish{\ttfamily ****}

\end{minipage}
\hfill
\begin{minipage}[t]{0.48\linewidth}

\textbf{\color{darkTeal}Space সহ}

\begin{english}
\ttfamily
for i in range(4):\\
\hspace*{4mm}print("*", end=" ")
\end{english}

\medskip

\textbf{Output:}
\textenglish{\ttfamily * * * *}

\end{minipage}

\medskip

দুইটিই সঠিক। প্রয়োজনীয় Output Design অনুযায়ী ব্যবহার করা হয়।

\end{codeoutputbox}

%================================================
% while Pattern
%================================================

\booksubheading{\textenglish{\ttfamily while} Loop দিয়েও Pattern তৈরি করা যায়}

Pattern সাধারণত
\textenglish{\ttfamily for}
Loop দিয়ে লেখা সহজ। তবে
\textenglish{\ttfamily while}
Loop দিয়েও একই Logic তৈরি করা যায়।

\begin{codeoutputbox}{\textenglish{\ttfamily while} দিয়ে Increasing Pattern}

\noindent
\begin{minipage}[t]{0.57\linewidth}
\codeheading

\begin{english}
\ttfamily
i = 1\\
while i <= 3:\\
\hspace*{4mm}j = 1\\
\hspace*{4mm}while j <= i:\\
\hspace*{8mm}print("*", end=" ")\\
\hspace*{8mm}j += 1\\
\hspace*{4mm}print()\\
\hspace*{4mm}i += 1
\end{english}

\end{minipage}
\hfill
\begin{minipage}[t]{0.39\linewidth}
\outputheading

\begin{english}
\ttfamily
*\\
* *\\
* * *
\end{english}

\medskip

Inner Counter
\textenglish{\ttfamily j}
প্রতিটি Outer Iteration-এ আবার
\textenglish{\ttfamily j = 1}
থেকে শুরু হয়।

\end{minipage}

\end{codeoutputbox}

\begin{rememberbox}

Pattern-এর জন্য সাধারণত
\textenglish{\ttfamily for}
Loop বেশি সহজ ও সংক্ষিপ্ত।
তবে
\textenglish{\ttfamily while}
Loop দিয়েও Pattern তৈরি করা যায়।

\end{rememberbox}

%================================================
% Same Loop, Different Print
%================================================

\booksubheading{একই Loop Logic, ভিন্ন Print Expression}

\begin{tabularx}
{\linewidth}
{
|>{\bfseries\RaggedRight\arraybackslash}p{0.39\linewidth}
|Y|
}

\hline
\rowcolor{softBlue}

\textbf{Code}
&
\textbf{ফলাফল}
\\
\hline

\textenglish{\ttfamily print("*", end=" ")}
&
Star Pattern
\\
\hline

\textenglish{\ttfamily print(i, end=" ")}
&
Row Number Pattern
\\
\hline

\textenglish{\ttfamily print(j, end=" ")}
&
Column Number Pattern
\\
\hline

\textenglish{\ttfamily print(num, end=" ")}
&
Continuous Number Pattern
\\
\hline

\end{tabularx}

%================================================
% Summary
%================================================

\booksubheading{গুরুত্বপূর্ণ Pattern Summary}

\begin{tabularx}
{\linewidth}
{
|>{\bfseries\RaggedRight\arraybackslash}p{0.34\linewidth}
|Y|
}

\hline
\rowcolor{softBlue}

\textbf{Pattern}
&
\textbf{Inner Loop / মূল Logic}
\\
\hline

Rectangle / Square
&
প্রতি Row-তে নির্দিষ্ট সংখ্যক Iteration।
\\
\hline

Increasing
&
\textenglish{\ttfamily range(i)}
বা Row-এর সঙ্গে Item সংখ্যা বৃদ্ধি।
\\
\hline

Decreasing
&
\textenglish{\ttfamily range(n-i+1)}
বা Reverse Count।
\\
\hline

Row Number
&
\textenglish{\ttfamily print(i)}
\\
\hline

Column Number
&
\textenglish{\ttfamily print(j)}
\\
\hline

Continuous Number
&
আলাদা Counter।
\\
\hline

\end{tabularx}

%================================================
% Common Mistakes
%================================================

\booksubheading{সাধারণ ভুল}

\begin{rememberbox}

\begin{itemize}

\item Outer Loop ও Inner Loop-এর কাজ গুলিয়ে ফেলা।

\item একই Row-তে Output রাখার জন্য
\textenglish{\ttfamily end}
ব্যবহার না করা।

\item Row পরিবর্তনের
\textenglish{\ttfamily print()}
ভুল Indentation-এ রাখা।

\item Increasing Pattern-এ Inner Loop-এর Range ভুল লেখা।

\item Decreasing Pattern-এর Item সংখ্যা ভুল হিসাব করা।

\item
\textenglish{\ttfamily i}
ও
\textenglish{\ttfamily j}
-এর কাজ গুলিয়ে ফেলা।

\item
\textenglish{\ttfamily range()}
-এর Stop Value Excluded—এটি ভুলে যাওয়া।

\item Nested
\textenglish{\ttfamily while}
Loop-এ Inner Counter পুনরায় Initialize না করা।

\item Star-এর জায়গায়
\textenglish{\ttfamily i}
বা
\textenglish{\ttfamily j}
ভুলভাবে Print করা।

\end{itemize}

\end{rememberbox}

%================================================
% Final Reminder
%================================================

\begin{rememberbox}

\textbf{\color{deepOrange}মনে রাখবে}

\begin{tabularx}
{\linewidth}
{
>{\bfseries\RaggedRight\arraybackslash}p{0.24\linewidth}
Y
>{\bfseries\RaggedRight\arraybackslash}p{0.24\linewidth}
Y
}

Outer Loop
&
Row
&
Inner Loop
&
প্রতি Row-এর Item
\\[0.5mm]

\textenglish{\ttfamily end=" "}
&
একই Line-এ Print
&
\textenglish{\ttfamily print()}
&
নতুন Row
\\[0.5mm]

Increasing
&
Item সংখ্যা বাড়ে
&
Decreasing
&
Item সংখ্যা কমে
\\[0.5mm]

\textenglish{\ttfamily print(i)}
&
Row Value
&
\textenglish{\ttfamily print(j)}
&
Column Value
\\[0.5mm]

Rectangle
&
Fixed Item Count
&
Continuous Number
&
আলাদা Counter
\\[0.5mm]

Nested Loop
&
Outer-এর প্রতি Iteration-এ Inner সম্পূর্ণ চলে
&
Pattern
&
Row + Column Logic
\\

\end{tabularx}

\medskip

\begin{center}

\textbf{Pattern Solve করার ৩ ধাপ}

\medskip

\textbf{কত Row?}
\(\longrightarrow\)
\textbf{Outer Loop}

\qquad

\textbf{প্রতি Row-তে কত Item?}
\(\longrightarrow\)
\textbf{Inner Loop}

\medskip

\textbf{কী Print হবে?}
\(\longrightarrow\)
\textenglish{\ttfamily * / i / j / counter}

\end{center}

\end{rememberbox}

\small

%================================================
% 5.43 Data Structure
%================================================

